\begin{center}
    \large
    
    \begin{singlespace}
        \textbf{\englishTitle{}} \\[0.5cm]
    \end{singlespace}
    
    \begin{singlespace}
        Student : \studentEnName{}  \hspace{1.0cm} 
        % 兩個指導教授要寫 Advisors
        \ifdefined\advisorCnNameB
            Advisors: Dr.\, \advisorEnName \\
            \hspace{6.6cm} Dr.\, \advisorEnNameB  \\
        \else
            Advisor: Dr.\, \advisorEnName \\
        \fi
    \end{singlespace}
    
    \vspace{0.5cm}
    \begin{singlespace}
        \NameofDepartmentInstituteEN\\
        National Yang Ming Chiao Tung University\\[0.2cm]
    \end{singlespace}
    
    \textbf{Abstract} \\[0.5cm]

\end{center}

\normalsize

This study investigates how model architecture, graph topology, and training-window length jointly determine cross-sectional return ranking and portfolio performance in the Taiwan equity market. Graph attention networks (GAT) and their neutralization-augmented variants are evaluated within a systematic ablation framework, using IC/ICIR-based signal quality and portfolio-level risk-adjusted return as dual criteria.

Using daily data from the Taiwan Economic Journal (TEJ), we construct 56 technical and valuation features and define the prediction target as the cross-sectionally demeaned five-day forward return. The experiment includes an industry graph based on intra-industry complete connections and a universe graph based on all-to-all market connections. Nine model configurations are compared: Linear regression, XGBoost, LSTM, and DNN as non-graph baselines, plus five GAT variants spanning different graph topology and neutralization designs, evaluated under short, medium, and long chronological training windows. Portfolio validation ranks stocks every five trading days, takes equal-weight long positions in the top decile, and benchmarks against the 0050 ETF.

Training-window length emerges as the dominant driver of test-set usability: short-window models yield positive Daily IC but produce near-zero or negative annualized portfolio returns, whereas medium and long windows sustain stable positive Sharpe ratios. Regarding graph topology, constraining GAT edges to intra-industry connections outperforms the all-to-all universe graph in both Daily IC and Sharpe ratio under medium and long windows, indicating that economically grounded local links better preserve stock-level cross-sectional dispersion. Augmenting GAT with industry-level representation neutralization achieves the highest long-window Daily IC (0.0590) and ICIR (0.786); further incorporating market-level neutralization yields the highest Sharpe ratio (1.849) and annualized return (42.02\%), revealing a trade-off between signal stability and absolute portfolio return. XGBoost reaches the highest long-window cumulative return (143.08\%) but exhibits a severe train--test IC gap (0.7824 versus 0.0205), indicating substantial overfitting risk. Overall, DNN without graph structure already delivers competitive cross-sectional ranking performance; the incremental contribution of graph attention mechanisms is best characterised as improvements in signal stability and topology interpretability under long training windows rather than uniform dominance over non-graph baselines.

% 5-7 Keywords (English)
\textbf{Keywords: Graph Attention Network, Cross-Sectional Return Prediction, Information Coefficient, Graph Topology Design, Taiwan Equities.}
